\documentclass[12pt]{article}
\usepackage[utf8]{inputenc}
\usepackage{amsmath, amssymb, amsfonts, mathtools, array}
\usepackage{geometry}
\geometry{a4paper, margin=1in}
\title{Student Answer Script Transcription}
\author{Roll Number: 240122044}
\date{\today}

\begin{document}
\maketitle

The student's handwritten notes indicate that the solution spaces of \(E\) and \(E'\) have the same RREF (Reduced Row Echelon Form). Therefore, using Elementary Row Operations (EROs), we can obtain both \(E\) and \(E'\) from their respective RREFs. Since \(E\) and \(E'\) are interchangable to each other, they are equivalent to each other.

\end{document}